This section describes how the webamc command-line tool can be used to
prepare, compile and submit MCQ that can be loaded via the web
interface. It assumes that you have already installed the webamc
package, following the instructions in Section~\ref{sec:install}.
%%%%%%%%%%%%%%%%%%%%
\subsection{Structuring tex files for webamc}
\label{subsec:submission/structuring}
MCQ submitted to the web interface may be of two forms:
\begin{itemize}
\item either a list of questions ;
\item or a list of exercises, each exercise being a list of questions.
\end{itemize}
From now on, the term {\em component} will be used to denote either an
MCQ, or an exercise, or a question. In addition to all components that
are part of it, an MCQ or an exercise may have a statement,
explaining, e.g., the context of all its components. This is the
purpose of files \file{\_mcq.tex} and \file{\_exercise.tex} to contain
this statements.

The following rules must be respected so that webamc can correctly
compile your code:
\begin{itemize}
\item The content of an MCQ or an exercise must be placed in a single
  directory.
\item To be recognised as an MCQ directory, the directory must contain
  an \file{\_mcq.tex} file containing the MCQ statement even if it is
  empty.
\item Likewise, to be recognised as an exercise directory, the
  directory must contain an \file{\_exercise.tex} file containing the
  exercise statement even if it is empty.
\item Questions must be placed in files of which the name is prefixed
  by \file{qst-}.
\end{itemize}

This is, for example, the directory structure of an MCQ containing
five question files:
\begin{mdframed}[style=defaultstyle]
\dirtree{%
.1 some\_mcq.
.2 \_header.tex.
.2 \_mcq.tex.
.2 qst-aaa.tex.
.2 qst-bbb.tex.
.2 qst-ccc.tex.
.2 qst-ddd.tex.
.2 qst-eee.tex.
}
\end{mdframed}

This is another example of the directory structure of a two exercises
MCQ (the first one with 2 question files and the second one with three
question files):
\begin{mdframed}[style=defaultstyle]
\dirtree{%
.1 some\_other\_mcq.
.2 \_header.tex.
.2 \_mcq.tex.
.2 exe-fst.
.3 \_exercise.tex.
.3 qst-fst-1.tex.
.3 qst-fst-2.tex.
.2 exe-snd.
.3 \_exercise.tex.
.3 qst-snd-1.tex.
.3 qst-snd-2.tex.
.3 qst-snd-3.tex.
}
\end{mdframed}

More complex examples may be found in the examples \file{examples/mcq}
directory.

You may also submit directories that do not contain any MCQ but only
question files. This is useful if you, for example, wish to record
some questions in the database before creating question packs (see
Section~\ref{subsec:submission/json-file-format}) that may contain
these question.

Note that all the file names listed above (\file{\_header.tex},
\file{\_exercise.tex}, \ldots) are only default names and may be
redefined through configuration files (see
Section~\ref{subsec:submission/configuration} for further details).
%%%%%%%%%%
\subsubsection{Header tex file - \file{\_header.tex}}
An \file{\_header.tex} file may contain latex code that will be placed
in the preamble of the latex document generated on-the-fly for each
component of the directory. Typically, you will write in this file all
\verb+\usepackage{...}+ instructions required to compile latex files
in this directory. A directory inherits from any \file{\_header.tex}
contained in any of its ancestor directory. For example, let's
consider this directory structure:
\begin{mdframed}[style=defaultstyle]
\dirtree{%
.1 some\_dir.
.2 \_header.tex~~~~~(content: \texttt{\textbackslash{}usepackage\{packa\}}).
.2 another\_dir.
.3 \_header.tex~~~~~(content: \texttt{\textbackslash{}usepackage\{packb\}}).
.3 \_exercise.tex~~~~~(content: This is a new exercise).
}
\end{mdframed}
When compiling the exercise file \file{\_exercise.tex}, the latex
document compiled will contain the following code:
\begin{tex}
\documentclass{article}
\def\WEBAMC{1}
\usepackage{packa}
\usepackage{packb}
\begin{document}
This is a new exercise
\end{document}
\end{tex}
%%%%%%%%%%
\subsubsection{Exercise tex file - \file{\_exercise.tex}}
An \file{\_exercise.tex} contains the tex code of an exercise
preamble, that is, the exercise statement required for the
understanding of the questions enclosed in this exercise. The
\name{CODE} attribute must be defined in this file (see
Section~\ref{subsec:submission/metadata}).
%%%%%%%%%%
\subsubsection{MCQ tex file - \file{\_mcq.tex}}
An \file{\_mcq.tex} contains the tex code of an MCQ preamble, that is
the MCQ statement required for the understanding of the components
(questions or exercises) enclosed in this MCQ. The \name{CODE} and
\name{TITLE} attributes must be defined in this file (see
Section~\ref{subsec:submission/metadata}).
%%%%%%%%%%
\subsubsection{Webamc metadata file - \file{\_webamc.tex}}
The content of this file is never compiled by webamc, but it may be
used to set some attributes for all the components enclosed in the
directory. For instance, if an MCQ directory contains a
\file{\_webamc.tex} file with the single line below:
\begin{tex}
%webamc RND=1
\end{tex}
then this means that the content of all components of the MCQ (i.e.,
its exercises, the questions within its exercises, and the answers of
all its questions) will be organised randomly.

As for \file{\_header.tex} files, a directory inherits from any
\file{\_webamc.tex} contained in any of its ancestor directory.
%%%%%%%%%%
\subsubsection{Question pack file - \file{\_pack.json}}
This is a JSON file containing the description of a pack of questions
that can be used to form an exercise or an MCQ. Please consult this
Section~\ref{subsec:submission/json-file-format} for further help on
the JSON file format.
%%%%%%%%%%
\subsubsection{Question file - \file{qst-XXX.tex}}
Any tex file which name starts with prefix \file{qst-} is considered
by webamc as a question file, that is, a file containing questions. A
question file may contain any tex code, but only the code enclosed
within questions will be compiled by webamc, e.g., tex code between
questions will be ignored.
%%%%%%%%%%%%%%%%%%%%
\subsection{Basics on compilation}
Right now, the compilation of the tex files that are part of your MCQ
is done on the client side using the \name{webamc} command line
tool. This has the advantage to alleviate the server but, since tex
files are compiled once before submitting the MCQ to the web
interface, this prevents MCQ from being dynamic.

When analysing your tex files, each question and each answer of your
MCQ will be compiled separately in order to a generate a separate png
file. For example, let's say your MCQ contains the following question:
\begin{tex}
\begin{questionmult}{math-sqrt}
    $\sqrt{81}$ = ?
    \begin{reponses}
    \bonne{9}
    \bonne{$3^2$}
    \mauvaise{19}
    \mauvaise{6}
    \mauvaise{8}
    \end{reponses}
\end{questionmult}
\end{tex}

Then webamc will on-the-fly generate and compile 6 tex documents (1
for the question statement and 1 for each answer) to produce 6 png
files. As a consequence, question and answer statements must be fully
context independent (e.g., with no \verb+\label{...}+ in it).

When compiling a latex file, webamc will define in the document's
preamble a macro called \verb+WEBAMC+ as follows:
\begin{tex}
\def\WEBAMC{1}
\end{tex}

This may be useful if you want to write latex code specific for the
web version of your MCQ, for example:
\begin{tex}
\ifdefined\WEBAMC
(Nota: This is the web version of ...)
\fi
\end{tex}

It is also very important to note that during compilation, all files
are compiled from the directory in which the \file{\_mcq.tex} file is
located. Therefore all paths that may appear in your question files
(e.g., if you include graphics in your questions) must be relative to
that directory.

However, if you compile a directory that does not contain an
\file{\_mcq.tex} file but only, e.g., question files, each question
file is then compiled from the directory it is located in.
%%%%%%%%%%%%%%%%%%%%
\subsection{Webamc metadata}
\label{subsec:submission/metadata}
Several component attributes can be set directly in tex files with
specific comments. Webamc will consider as an attribute definition any
comment that starts at column one and that follows this pattern:
\begin{tex}
%webamc ATTRIBUTE_NAME=ATTRIBUTE_VALUE
\end{tex}
Note that once loaded in the web interface all these attributes
(except the \name{CODE}) can be modified through the web interface.
%%%%%%%%%%
\subsubsection{Attribute list}
%%%%%
\paragraph{Attribute \name{CODE}}
This defines the code of the component. Codes must be unique and not
null. Therefore this attribute is mandatory in exercise tex files
(\file{\_exercise.tex}) and MCQ tex files (\file{\_mcq.tex}). For a
question, the code is directly extracted from the tex code (i.e.,
\verb+id-qst+ in \verb+\begin{question}{id-qst}+).
%%%%%
\paragraph{Attribute \name{DIFFICULTY}}
This defines the difficulty of the component. It is an integer ranging
from 1 to 5. This is useful if you want to create question packs based
on their difficulty (see this
Section~\ref{subsec:submission/json-file-format}).
%%%%%
\paragraph{Attribute \name{RND}}
This defines whether or not the component content (e.g., questions of
an exercise or answers of a question) is randomized. Accepted values
are \verb+0+ and \verb+1+. Default value is \verb+0+.
%%%%%
\paragraph{Attribute \name{STANDALONE}}
Only valid for questions (for now), this defines whether or not the
question is a standalone question that can be selected in question
packs. Non standalone questions are typically context dependent
questions (e.g., questions based on an exercise statement).  Accepted
values are \verb+0+ and \verb+1+. Questions that are not part of an
exercise or an MCQ are always considered as standalone
questions. Otherwise, the default value is \verb+0+.
%%%%%
\paragraph{Attribute \name{TAG}}
This adds a tag to the component. Several tags can be added through
repeated assignments to that attribute. This is useful if you want to
create question packs based on their tags (see this
Section~\ref{subsec:submission/json-file-format}).
%%%%%
\paragraph{Attribute \name{TITLE}}
This defines the title of the component that will appear in the web
interface. Note that any tex macro appearing in it will be ignored,
i.e., not compiled, and will appear as such in the web interface. This
attribute is mandatory in MCQ tex files (\file{\_mcq.tex}).
%%%%%
\paragraph{Attribute \name{VISIBLE}}
This defines the visibility of the component. A non visible component
will not appear in the web interface, except for the user that
submitted it. Default value is \verb+1+.
%%%%%%%%%%
\subsubsection{Attribute validity}
Some attributes are mandatory in some files whereas some others may be
meaningless in others. The table below synthesises attribute
availability in the different tex files analysed by webamc. In this
table:
\begin{itemize}
\item 0 means that, if present in the file, the attribute will be
  ignored ;
\item 1 means that the attribute will be considered in the file
  although it is not mandatory ;
\item and 2 means that the attribute is mandatory in the file.
\end{itemize}

\begin{center}
  \begin{tabular}{lrrrrr}
    & \file{\_exercise.tex} & \file{\_header.tex}
    & \file{\_mcq.tex} & \file{\_webamc.tex} & \file{qst-XXX.tex}\\
    \hline
    \name{CODE}       & 2 & 0 & 2 & 0 & 0\\
    \name{DIFFICULTY} & 1 & 0 & 1 & 1 & 1\\
    \name{RND}        & 1 & 0 & 1 & 1 & 1\\
    \name{STANDALONE} & 0 & 0 & 0 & 1 & 1\\
    \name{TAG}        & 1 & 0 & 1 & 1 & 1\\
    \name{TITLE}      & 1 & 0 & 2 & 0 & 0\\
    \name{VISIBLE}    & 1 & 0 & 1 & 1 & 1\\
  \end{tabular}
\end{center}
%%%%%%%%%%%%%%%%%%%%
\subsection{Component compilation and submission}
The \command{webamc compile} command takes two arguments:
\begin{itemize}
\item a directory containing some component(s) to compile ;
\item and the name of the zip archive (without the \file{.zip} suffix)
  that you will get in output.
\end{itemize}

You will find several examples in the \file{examples} directory of the
package. Note that this process may take a while (typically a few
minutes for large MCQ) since, as said above, webamc will compile each
component and each answer separately. The zip archive you get in output
can be submitted through the web interface.
%%%%%%%%%%%%%%%%%%%%
\subsection{JSON file format for question packs}
\label{subsec:submission/json-file-format}
When present in an exercise or an MCQ directory, a \file{\_pack.json}
file specifies that this one is composed of a set of standalone
questions (i.e., questions for which the \name{STANDALONE} attribute
has been set to \verb+1+). The purpose of the \file{\_pack.json} is to
specify how to pick these questions. It consists basically of a
sequence of operations successively applied. Available operations are:
\begin{itemize}
\item \verb+all+: select all available questions
\item \verb+sort+: sort questions (by \verb+difficulty+ or by
  \verb+code+)
\item \verb+shuf+: select questions randomly
\item \verb+head+: select the first questions from those selected so far
\item \verb+with-code+: select questions based on their codes
\item \verb+with-difficulty+: select questions based on their
  difficulty
\item \verb+with-tag+: select questions based on their tags
\end{itemize}

An operation is specified with a dictionary having the following keys:
\begin{itemize}
\item \verb+op+: the operation realised, among those listed above
\item \verb+arg+: for some operations, specify an argument (e.g., the
  number of questions selected for operation \verb+head+)
\item \verb+rev+: for \verb+with-XXX+ operations, operate a negation
  during the selection (e.g., to select questions that do {\em not}
  have a specific tag)
\item \verb+content+: questions on which this operation is realised,
  also specified using an operation dictionary. By default all
  available questions are taken.
\end{itemize}
  
Here are a few examples:

\begin{itemize}
\item select all available questions from the database:
  \lstinputlisting[language=json]{../examples/mcq/packs/all/_pack.json}

\item select all available questions from the database and order them
  randomly:
  \lstinputlisting[language=json]{../examples/mcq/packs/shuf/_pack.json}

\item randomly pick 5 questions:
  \lstinputlisting[language=json]{../examples/mcq/packs/pick/_pack.json}

\item take all questions with tag \name{réseaux} and with a difficulty
  of 4 or 5 and shuffle them:
  \lstinputlisting[language=json]{../examples/mcq/packs/tags/_pack.json}

\item take all questions that do not have the \name{linux} tag:
  \lstinputlisting[language=json]{../examples/mcq/packs/tags-rev/_pack.json}

\item take all questions having the \name{réseaux} tag and sort them
  by \verb+difficulty+:
  \lstinputlisting[language=json]{../examples/mcq/packs/sort/_pack.json}
\end{itemize}

It is also possible to realise the union of several selections using
lists of operations. For example, this specifies a question pack with
all questions having the \name{linux} tag or with a difficulty level
strictly less than 3 ordered randomly.
\lstinputlisting[language=json]{../examples/mcq/packs/union/_pack.json}

Last, note that the web interface also provides a form to directly
create question packs rather than specifying them with a JSON file.
%%%%%%%%%%%%%%%%%%%%
\subsection{Compilation configuration}
\label{subsec:submission/configuration}
Several configuration parameters alter the compilation process.
Therefore after installation you may wish to check your configuration
file, although the default values should work
fine. Section~\ref{sec:config} also gives the default configuration.

These parameters configure the way the tex-to-pdf compiler is invoked:
\begin{itemize}
\item \texttt{tex2pdf\_exe}: name of the executable (that can be found
  in your \name{PATH} environment variable), or absolute path of the
  executable that will be used to compile tex files to pdf files
\item \texttt{tex2pdf\_exe\_args}: argument list given to this
  executable. The two following strings {\em must} appear in this
  list:
  \begin{itemize}
  \item \texttt{\{out\_dir\}}: placeholder for the directory name in
    which the tex-to-pdf compiler will store output file
  \item \texttt{\{tex\_file\}}: placeholder for the input tex file
  \end{itemize}
\item \texttt{log\_file}: path of the log file in which webamc will
  store compilation results
\end{itemize}

You may also change the name of the default file names that webamc
will handle (as described in
Section~\ref{subsec:submission/structuring}) by changing these
configuration parameters:
\begin{itemize}
\item \texttt{json\_file\_pack}
\item \texttt{tex\_file\_exercise}
\item \texttt{tex\_file\_header}
\item \texttt{tex\_file\_mcq}
\item \texttt{tex\_file\_question\_prefix}
\item \texttt{tex\_file\_webamc}
\end{itemize}

You can also specify which tex macros or environments will be handled
by webamc with these configuration parameters:
\begin{itemize}
\item \texttt{tex\_envs\_choices}
\item \texttt{tex\_envs\_question}
\item \texttt{tex\_envs\_question\_mult}
\item \texttt{tex\_macros\_choice}
\item \texttt{tex\_macros\_choice\_correct}
\item \texttt{tex\_macros\_lastchoices}
\end{itemize}
